% Generated draft: requires mathematical and source review.
\documentclass[11pt]{article}
\usepackage[margin=25mm]{geometry}
\usepackage{fontspec}
\setmainfont{Latin Modern Roman}
\usepackage{amsmath,amssymb,mathtools,unicode-math}
\setmathfont{Latin Modern Math}
\usepackage{enumitem,needspace,xurl,hyperref}
\setlength{\parindent}{0pt}
\setlength{\parskip}{5pt}
\setlength{\emergencystretch}{3em}
\newcommand{\N}{\mathbb{N}}
\newcommand{\Z}{\mathbb{Z}}
\newcommand{\Q}{\mathbb{Q}}
\newcommand{\R}{\mathbb{R}}
\newcommand{\C}{\mathbb{C}}
\newcommand{\F}{\mathbb{F}}
\newcommand{\A}{\mathbb{A}}
\newcommand{\PP}{\mathbb{P}}
\newcommand{\E}{\mathbb{E}}
\newcommand{\eps}{\varepsilon}
\newcommand{\dd}{\,\mathrm{d}}
\begin{document}
\section*{UnsolvedMath ID: 800416}
\textbf{Source problem code:} catalogue entry
\subsection{Definitions and mathematical statement}
For a self-map \(g\) and a point \(Q\), define the forward orbit by
\[
\mathcal O_g(Q)=\{g^n(Q):n\in\Z_{\geq 0}\},\qquad g^0=\operatorname{id}.
\]
Its cardinality counts distinct points. An orbit is Zariski dense if its Zariski closure is the whole ambient variety.

Reduction is understood at primes where the indicated map and point have good reduction. In the affine case, exclude primes dividing denominators of the coordinates of \(P\) or coefficients of \(f\) or its polynomial inverse; then \(\widetilde f\) is an automorphism over \(\F_p\). In the projective case, require that reduction of \(f\) remains a morphism of degree \(d\). These exclusions remove only finitely many primes. Write \(\widetilde P\) for the reduced point and \(\widetilde f\) for the reduced map.

The supplied source says that certain sets have ``density \(1\)'' without specifying the notion of density. Under the usual interpretation as relative natural density among the primes, this means
\[
\lim_{x\to\infty}
\frac{\#\{p\leq x:p\text{ is prime and }p\in S\}}
{\pi(x)}=1,
\qquad
\pi(x)=\#\{p\leq x:p\text{ is prime}\}.
\]
The formulation below uses this interpretation explicitly; the intended density convention requires source confirmation. The source also leaves the range of the dimension parameter \(N\) implicit.

\begin{enumerate}
\item[(a)] Let \(f:\PP^N\to\PP^N\) be a morphism defined over \(\Q\), of degree \(d\geq2\), and let \(P\in\PP^N(\Q)\) have Zariski dense forward orbit. The conjecture asserts that for every real \(\epsilon>0\), the set
\[
S_{f,P,\epsilon}^{\mathrm{proj}}
=
\left\{p\text{ prime of good reduction}:
\#\mathcal O_{\widetilde f}(\widetilde P)
\geq p^{(1/2-\epsilon)N}\right\}
\]
has density \(1\) among the primes, with the density qualification stated above.

\item[(b)] Let \(f:\A^N\to\A^N\) be a polynomial automorphism defined over \(\Q\), meaning that \(f\) is given by polynomials over \(\Q\) and has a polynomial inverse over \(\Q\). Let \(P\in\A^N(\Q)\) have Zariski dense forward orbit. The conjecture asserts that for every real \(\epsilon>0\), the set
\[
S_{f,P,\epsilon}^{\mathrm{aff}}
=
\left\{p\text{ prime of good reduction}:
\#\mathcal O_{\widetilde f}(\widetilde P)
\geq p^{(1-\epsilon)N}\right\}
\]
has density \(1\) among the primes, with the same density qualification.
\end{enumerate}
In each assertion, \(f\), \(P\), and \(\epsilon\) are fixed before taking the density limit.
\subsection{Short English statement}
The supplied conjecture predicts that a rational point with Zariski dense orbit has a large orbit after reduction modulo almost every prime: at least \(p^{(1/2-\epsilon)N}\) points for a projective morphism of degree at least \(2\), and at least \(p^{(1-\epsilon)N}\) points for an affine polynomial automorphism, for every \(\epsilon>0\). Here ``almost every'' uses the explicitly qualified interpretation of density among primes. The registry labels this problem open; that status has not been independently verified.
\subsection{Sources}
\begin{itemize}
\item Local UnsolvedMath registry, permanent ID \(800416\), exact source problem code \texttt{catalogue entry}, title ``Open Problems in Arithmetic Dynamics --- Conjecture 39.'' The supplied record has upstream revision \texttt{d371ee917ac9647fbb06cd6fc529deec6f3dda82} and status \texttt{open}.
\item Benedetto, DeMarco, Ingram, Jones et al., \emph{Current Trends and Open Problems in Arithmetic Dynamics} (2018), Conjecture 39. Supplied source URL: \url{https://arxiv.org/abs/1806.04980}. The registry attributes the proposal to these authors and gives the proposed year as 2018; separate original attribution is not supplied.
\item Supplied provenance: source extraction method \texttt{arxiv\_latex}; recorded access date 2026-10-16. This date and provenance are supplied source metadata, not an independent source visit or verification by this draft. The registry record was created and updated on 2026-10-16. No separate dated literature assessment is supplied.
\end{itemize}
\end{document}
